\section{Radial spatial geometry and aligned profiles}\label{sec:radial}

The quiet criterion supplies the testing regions, where $G\ge V$.
A center change requires a different arrangement: $V$ must be close to
$\lambda G$ with a large scalar $\lambda$. This section constructs radial
spatial metrics that can pass between these arrangements and proves the
curvature estimate for the aligned regime. Here \emph{aligned} means
that their metric jets through order two are close to proportional,
with the scalar frozen at the evaluation point.

All norms of covariant jets in this section are multilinear operator norms in the stated fixed affine coordinates. The norm of a Hermitian contravariant matrix is its Frobenius norm. Constants called absolute may depend on a fixed upper bound for the first two derivatives of $\kappa$ in $\log Y$, but not on the spatial complex dimension. This convention is essential when choosing the dimension only after the geometric constants.

\subsection{The radial clocks}

Write $t=|s|^2$, $\nu=\log t$, and let $B$ and $\Psi$ be radial potentials. Put
\begin{equation}\label{radial:definitions}
 \begin{gathered}
 a=B'(t),\qquad a_\psi=\Psi'(t)=a/\lambda,\qquad
 u=\log a,\\
 Y=\frac{du}{d\nu},\qquad
 H=\frac{Y}{t a_\psi},\qquad
 p=\frac{d\log a_\psi}{du},\qquad
 \kappa=\frac{dY/d\nu}{Y(1+Y)}.
 \end{gathered}
\end{equation}
Here a prime on a radial potential means differentiation in $t$. We use $u$ for this radial variable in this section and $u_s$ in Section~\ref{sec:assembly}. We write $V=\partial\bar\partial B$ and $G_0=\partial\bar\partial\Psi$ for their Hermitian matrices.

These quantities separate the controls needed in the construction.
The ratio $\lambda$ compares the transverse coefficients of $V$ and
$G_0$, while $H/\lambda=Y/(ta)$ is the common scale of the spatial
curvature components computed below. The evolution equations will give
$d\Psi/du=1/H$ and $d\nu/du=1/Y$: large $H$ limits the change of the
radial twist $\Psi$ during growth, and large $Y$ limits the spatial radius consumed.
We use the first feature to prepare a center change and the second to
keep each accelerated episode inside a finite radial interval. Capped
growth then lets the spatial radius continue between successive tests.

\begin{lemma}[Radial evolution]\label{radial:odes}
Suppose $a,a_\psi,Y>0$, $\lambda\ge1$, and
\begin{equation}\label{radial:kappa-assumptions}
 \frac1{1+Y}\le\kappa\le1.1,\qquad
 \frac{d\kappa}{d\log Y}\ge-\kappa.
\end{equation}
Then
\begin{align}\label{radial:clock-equations}
 \frac{d\nu}{du}&=\frac1Y,&
 \frac{d\log Y}{du}&=\kappa(1+1/Y),&
 \frac{d\log\lambda}{du}&=1-p,\\
 \frac{d\Psi}{du}&=\frac1H,&
 \frac{d\log H}{du}&=\kappa(1+1/Y)-p-1/Y.\notag
\end{align}
If $Y\ge t$ initially, then $Y\ge t$ subsequently. With smooth bounded controls $p,\kappa$ on each finite $u$ interval, positive solutions have no blowup on that interval. A cap of the form $\kappa=C_\kappa/(1+Y)$, with fixed $C_\kappa\ge1$, has no blowup at finite $\nu$ for the prescribed bounded controls.

If $p=1$, then
\begin{equation}\label{radial:H-logY}
 \frac{d\log H}{d\log Y}=1-\frac1\kappa.
\end{equation}
In particular, $p=\kappa=1$ holds both $\lambda$ and $H$ constant. If $\kappa\ge1$ on an interval beginning at $u_0$, then the total increase of $\nu$ on that interval is at most $1/Y(u_0)$, even if its finite $u$ length is arbitrarily large.
\end{lemma}

\begin{proof}
The first three identities follow by changing the independent variable from $\nu$ to $u$. Since $d\Psi/d\nu=t a_\psi$ and $H=Y/(t a_\psi)$, the fourth follows as well. Differentiating $\log H=\log Y-\nu-\log a_\psi$ gives the fifth. The lower bound on $\kappa$ implies
\[
 \frac{d\log(Y/t)}{d\nu}=\kappa(1+Y)-1\ge0.
\]
In $u$ time, $dY/du=\kappa(1+Y)\le1.1(1+Y)$, so $Y$ remains finite and positive on each finite interval. The remaining variables are obtained by integrating the displayed equations there. In a cap, $dY/d\nu=C_\kappa Y$, which gives the asserted global $\nu$ continuation. Dividing the last equation of \eqref{radial:clock-equations}, with $p=1$, by the second proves \eqref{radial:H-logY}. Finally, $\kappa\ge1$ gives $Y(u)\ge Y(u_0)e^{u-u_0}$, and integration of $d\nu/du=1/Y$ gives the last assertion.
\end{proof}

The inequalities in \eqref{radial:kappa-assumptions} allow the smooth changes used later. For example, to decrease $\kappa$ between two fixed positive values, make the change over a sufficiently long fixed interval in $\log Y$; then $d\log\kappa/d\log Y\ge-1$. To pass from $\kappa=1$ to a cap, choose a smooth function $K(Y)$ which initially equals $1+Y$, has $0\le K'\le1$, and eventually is constant, and put $\kappa=K(Y)/(1+Y)$. The change can be made on a fixed increasing-factor interval of $Y$, so its first two $\log Y$ derivatives are bounded. Moreover $1\le K\le1+Y$ and
\[
 Y\kappa_Y+\kappa
 =\frac{K+Y(1+Y)K'}{(1+Y)^2}\ge0.
\]
Thus all inequalities persist, and $H$ is nonincreasing during this change when $p=1$.

\subsection{Curvature and coordinate bounds}

\begin{lemma}[Radial curvature]\label{radial:curvature}
Under \eqref{radial:kappa-assumptions}, the negative Hermitian curvature tensor $C^V=-R^V$ satisfies
\begin{equation}\label{radial:rank-reserve}
 C^V(T,T)\ge c\,\kappa\frac{H}{\lambda}|T|_V^2
 \qquad\text{if }T=T^*,\quad\operatorname{rank}T\le2,
\end{equation}
where $c>0$ is absolute. Here the subscript $V$ means that the matrix is expressed in $V$-unit axes.
\end{lemma}

\begin{proof}
At $s=(\sqrt t,0,\ldots,0)$ the metric has entries
\[
 V_{i\bar j}=a\delta_{ij}+a'\bar s_i s_j,
\]
and its radial and transverse eigenvalues are $a(1+Y)$ and $a$. The unitary symmetry in the transverse variables and the K\"ahler symmetries leave only a transverse block, a radial--transverse bisectional block, and a pure radial entry. After removing the common factor $Y/(ta)=H/\lambda$, their respective values in $V$-unit axes are
\begin{equation}\label{radial:curvature-components}
 \delta_{ij}\delta_{kl}+\delta_{il}\delta_{kj},\qquad
 \kappa\delta_{kl},\qquad
 D=\kappa\bigl(1+\kappa+Y\kappa_Y\bigr).
\end{equation}
For completeness, the transverse coefficient is $a'/a^2=Y/(ta)$. The radial--transverse coefficient is
\[
 \frac{dY/d\nu}{t a(1+Y)}=\kappa\frac{Y}{ta}.
\]
For the pure radial entry, differentiate the logarithm of the radial metric along the radial complex line. Since
\[
 \frac{d}{d\nu}\log\bigl(a(1+Y)\bigr)=Y(1+\kappa),
\]
its second $\nu$ derivative divided by $ta(1+Y)$ is exactly the last expression in \eqref{radial:curvature-components} times $Y/(ta)$. These calculations also follow directly by subtracting the cubic contractions in the K\"ahler curvature formula.

Write a Hermitian test matrix in radial/transverse blocks as
\[
 T=\begin{pmatrix}r&b^*\\ b&S\end{pmatrix}.
\]
The contraction after removing $Y/(ta)$ is
\[
 Dr^2+2\kappa r\operatorname{tr}S+2\kappa|b|^2
       +(\operatorname{tr}S)^2+|S|^2.
\]
The assumption on $\kappa_Y$ gives $D\ge\kappa$. Absorb the cross term with $0.9Dr^2$. Its cost is at most
\[
 \frac{\kappa^2}{0.9D}(\operatorname{tr}S)^2
 \le\frac{1.1}{0.9}(\operatorname{tr}S)^2.
\]
Since $\operatorname{rank}S\le2$, $|\operatorname{tr}S|\le\sqrt2|S|$. The transverse reserve is therefore at least $5|S|^2/9$, in addition to $0.1Dr^2$ and $2\kappa|b|^2$. As $\kappa\le1.1$, these bound an absolute positive multiple of $\kappa(r^2+2|b|^2+|S|^2)$. This proves \eqref{radial:rank-reserve}.
\end{proof}

\begin{lemma}[Scaled affine jets]\label{radial:scaled}
Assume \eqref{radial:kappa-assumptions} and $\lambda\ge1$.
Suppose $Y\ge1$, $Y\ge t$, the first two derivatives of $\kappa$ in $\log Y$ are bounded by a fixed constant, and
\begin{equation}\label{radial:p-close}
 |p-1|+\left|\frac{dp}{d\log Y}\right|
      +\left|\frac{d^2p}{d(\log Y)^2}\right|\le\delta,
 \qquad 0<\delta\le\tfrac1{10}.
\end{equation}
At a point use $G_0$-unit affine axes divided by $\sqrt H$. In these axes $G_0=H^{-1}I$ at the point. The transverse and radial Euclidean axis lengths are
\begin{equation}\label{radial:axis-lengths}
 \sqrt{\frac tY},\qquad
 \sqrt{\frac{t}{Y(1+pY)}}.
\end{equation}
In these fixed coordinates:
\begin{enumerate}
\item derivatives of $\Psi$ and $B$ in orders one through four are bounded by $C/H$ and $C\lambda/H$, respectively;
\item metric jets through order two of $V-\lambda G_0$, with $\lambda$ frozen at its value at the point, are bounded by $C\delta\lambda/H$;
\item derivatives of $u$ in orders one through four are bounded by $C$;
\item if a further potential has Euclidean derivatives in orders two through four bounded by $D_e$, its Hessian and first two metric jets here are bounded by $CD_e/(a_\psi H)$.
\end{enumerate}
All constants are independent of dimension.
\end{lemma}

\begin{proof}
The eigenvalues of $G_0$ are $a_\psi$ transversely and $a_\psi(1+pY)$ radially, giving \eqref{radial:axis-lengths}. In these axes $dt$ and $d^2t$ have operator norms at most $Ct/Y$, and higher derivatives of $t$ vanish. The first bound includes the fact that the only nonzero first derivative is radial; thus its axis length includes the extra factor $(1+pY)^{-1/2}$.

Set a dot to mean $d/d\nu$. We have
\[
 \dot Y=\kappa Y(1+Y)=O(Y^2),\qquad
 \ddot Y=O(Y^3),\qquad \dddot Y=O(Y^4).
\]
The last two estimates use, respectively, the first and second indicated derivatives of $\kappa$. Also $\dot p=O(\delta Y)$ and $\ddot p=O(\delta Y^2)$. Direct differentiation gives
\begin{align*}
 a'&=aY/t,\\
 a''&=a\bigl(Y^2-Y+\dot Y\bigr)/t^2,\\
 a'''&=a\bigl(Y^3-3Y^2+2Y+3(Y-1)\dot Y+\ddot Y\bigr)/t^3.
\end{align*}
Consequently $|a^{(j)}|\le Ca(Y/t)^j$ for $0\le j\le3$. The same formulas for $a_\psi$ use $pY$ in place of $Y$ and its derivatives. They imply the corresponding estimate for $a_\psi$, and comparison at the point, where $a=\lambda a_\psi$, improves their difference by the factor $\delta$.

The chain rule for a radial potential in at most four spatial slots is a finite sum of products of its $t$ derivatives and copies of $dt,d^2t$. For a term containing its $l$th $t$ derivative, the product of these latter factors is bounded by $C(t/Y)^l$. Since $B^{(l)}=a^{(l-1)}$, every such product is at most $Cat/Y=C\lambda/H$, and the same argument gives $C/H$ for $\Psi$. The improved difference estimate proves the metric-jet assertion. For $u$, its $t$ derivatives through order four are $O((Y/t)^j)$, by the displayed bounds on $Y$ and its derivatives; the same chain rule proves the third assertion. All sums here have a number of terms bounded in terms of the derivative order, not the dimension.

Finally, the largest axis length in \eqref{radial:axis-lengths} is $\sqrt{t/Y}\le1$. Applying two, three, or four such axes to the extra potential therefore gives at most $CD_e t/Y=CD_e/(a_\psi H)$.
\end{proof}

\subsection{Constant-center aligned estimates}

In the remainder of this section the fiber center and severity $E$ are constant, the fiber profile is $q=q_0(x,E)$ from Proposition~\ref{quiet:profile}, and no shear is added. We use the fiber notation
\[
 x=\log\tau,\quad P=\tau q,\quad r=q_x/q,\quad
 D_q=(q_x+q_{xx})/q.
\]
The elementary profile estimates of Lemma~\ref{quiet:scalar} give $q\ge1$, $r>0$, $D_q\ge r^2$, and, for constants allowed to depend on the fixed profile and $E$,
\begin{equation}\label{radial:profile-bounds}
 \begin{split}
 &D_q\ge c_E\tau\quad(0<\tau\le1),\qquad q\le C_E\quad(0<\tau\le1),\\
 &q/D_q\le C q^2,\qquad q\le C_E(1+\log\tau)\quad(\tau\ge1).
 \end{split}
\end{equation}
These bounds also follow directly from the positive logarithmic summand of $q_0$: for $\tau\ge1$ it supplies a positive lower bound for $q_x+q_{xx}$, whereas the other summand is bounded in $\tau$ at fixed $E$.

\begin{proposition}[Aligned profile]\label{radial:aligned}
There is an absolute $\delta_0>0$ with the following property. Fix $\sigma>0$ and take the severity $E$ sufficiently large in terms of $\sigma$ and the fixed profile. Assume the hypotheses of Lemma~\ref{radial:scaled}, with $\delta\le\delta_0$, and also
\[
 \tfrac12\le\kappa\le1.1,\qquad H\ge\sigma,
 \qquad G=G_0+G_e,\qquad
 |\partial^{[j]}G_e|\le\delta/H\quad(0\le j\le2)
\]
in its scaled affine axes. Here $G=\partial\bar\partial\psi$ and $\partial^{[j]}$ denotes any real or complex metric jet of order $j$, measured as a multilinear operator. Suppose either
\begin{equation}\label{radial:aligned-alternatives}
 \lambda=1,\qquad\text{or}\qquad
 H\ge M_E(1+\log\lambda)^2,
\end{equation}
where $M_E$ is sufficiently large depending only on the fixed profile and $E$. Then the constant-center fiber potential defines a positive metric with nonpositive real sectional curvature. Its spatial Schur complement satisfies $h\ge cV$.

More precisely, in scaled spatial axes and the tilted logarithmic fiber coordinate of Lemma~\ref{fiber:blocks}, the contraction of its negative curvature on each Hermitian rank-at-most-two matrix with blocks $(T,b,t_f)$ is bounded below by
\begin{equation}\label{radial:aligned-margin}
 c\left(A t_f^2+W(b,\bar b)
               +\frac{\lambda+\tau}{H}|T|^2
               +P|T|_G^2\right),
\end{equation}
where $c>0$ can be fixed independently of dimension and the stage parameters once the stated thresholds hold. The coefficient $A$ and the form $W$ are those of Lemma~\ref{fiber:blocks}.
\end{proposition}

\begin{proof}
Since $q$ has no spatial dependence, $\mathcal N=G$ and $h=V+\tau G$. Lemma~\ref{radial:scaled} shows that, through two metric jets, $h$ differs from the constant multiple $(\lambda+\tau)V/\lambda$ by at most $C\delta(\lambda+\tau)/H$. In particular
\[
 h\asymp\frac{\lambda+\tau}{H}I,\qquad
 \Gamma^h=O(1).
\]
Lemma~\ref{radial:curvature} gives a covariant curvature lower bound $c\lambda|T|^2/H$ for $V$ in these scaled axes. Multiplication of a metric by a positive constant multiplies its covariant curvature by that constant. The curvature difference is at most $C\delta(\lambda+\tau)|T|^2/H$: indeed its fourth-derivative term has this bound, and its cubic contractions have the same bound by the inverse estimate $|h^{-1}|\le CH/(\lambda+\tau)$ and the first-jet bounds. Lemma~\ref{fiber:perturbation} makes this calculation explicit. Reducing $\delta_0$ therefore gives
\begin{equation}\label{radial:h-margin}
 C^h(T,T)\ge c\frac{\lambda+\tau}{H}|T|^2
 \quad(\operatorname{rank}T\le2).
\end{equation}

Let $\theta=\lambda/(\lambda+\tau)$. The identity $\nabla^h(V+\tau G)=0$ implies
\[
 \nabla^h G=\theta\bigl(\nabla^hG-\lambda^{-1}\nabla^hV\bigr),
 \qquad |\nabla^hG|\le C\delta\theta/H.
\]
The curvature table gives $A=Pq^2D_q$, vanishing spin and $A_i$ terms, and $K=P\nabla^hG$. In $G$-unit normalization it gives
\begin{equation}\label{radial:W-normalized}
 \frac{W}{Pq}=rI+\mathcal Y,
 \qquad \mathcal Y=V_G(V_G+\tau I)^{-1},
 \qquad \mathcal Y\asymp\theta I.
\end{equation}
Here $V_G$ denotes $V$ expressed in $G$-unit axes; its eigenvalues are $(1+O(\delta))\lambda$. In particular $W>0$. Spending a fixed small portion of $W$ on the $K$ coupling costs at most
\[
 C\delta^2\frac{P\theta^2}{qH(r+\theta)}|T|^2
 \le C\delta^2\frac{\tau\theta}{H}|T|^2,
\]
which is an arbitrarily small portion of \eqref{radial:h-margin} after reducing $\delta_0$.

It remains to pay for $2t_f(W:T)$. Spending $0.98A t_f^2$ costs $(W:T)^2/(0.98A)$. First suppose the second alternative in \eqref{radial:aligned-alternatives} holds. Split the two terms in \eqref{radial:W-normalized} using $(a+b)^2\le1.1a^2+11b^2$. The $r$ term costs at most
\[
 \frac{1.1}{0.98}P(\operatorname{tr}_G T)^2.
\]
The horizontal curvature contains the positive term
\[
 P\bigl((\operatorname{tr}_G T)^2+|T|_G^2\bigr).
\]
As $|\operatorname{tr}_G T|\le\sqrt2|T|_G$, their difference retains more than $P|T|_G^2/2$. The remaining trace cost is at most
\begin{equation}\label{radial:trace-remainder}
 C\frac{P\theta^2}{D_qH^2}|T|^2.
\end{equation}
For $\tau\le1$, \eqref{radial:profile-bounds} bounds $P/D_q$ by $C_E$. For $\tau\ge1$, it bounds \eqref{radial:trace-remainder} by
\[
 \frac{C_E}{H^2}\tau\theta^2(1+\log\tau)^2|T|^2
 \le \frac{C_E(1+\log\lambda)^2}{H^2}
             (\lambda+\tau)|T|^2.
\]
For the last inequality, put $z=\tau/\lambda$ and use boundedness of
$z(1+z)^{-3}$ and $z(1+z)^{-3}(\log z)^2$ on $(0,\infty)$. Taking $M_E$ sufficiently large absorbs this cost into a small portion of \eqref{radial:h-margin}.

Now suppose $\lambda=1$. The profile estimates give, for sufficiently large $E$,
\begin{equation}\label{radial:lambda-one-slope}
 \|\mathcal Y\|\le r/18\qquad(\tau>2e^{-E}).
\end{equation}
To see the uniformity of this assertion, $\|\mathcal Y\|\le(1+C\delta)/(1+\tau)$. On $2e^{-E}\le\tau\le1$, the first summand of $q_0$ gives a logarithmic slope bounded below by a large fixed multiple of one. On $1\le\tau\le e^{3b_0E}$, it gives $r\ge0.9b_0(1+\tau)^{-1/16}$ after increasing $E$. Beyond this interval the logarithmic summand gives $r\ge c/q$, which exceeds $40/(1+\tau)$ by its linear-in-$\log\tau$ upper bound and the starting size $\tau\ge e^{3b_0E}$. These statements are unchanged by taking $\delta_0$ smaller.

By \eqref{radial:lambda-one-slope}, the trace expenditure outside the exceptional interval is at most
\[
 \frac{P}{0.98}
       \left(|\operatorname{tr}_GT|+\frac{\sqrt2}{18}|T|_G\right)^2.
\]
For $0\le s\le\sqrt2$, the ratio $(s+\sqrt2/18)^2/[0.98(1+s^2)]$ is strictly less than one, by an absolute amount. Thus a fixed portion of $P|T|_G^2$ remains. On $\tau\le2e^{-E}$, the explicit left-end profile gives $q\le C$, $r\le C$, and $D_q\ge c\tau e^E$, where constants may depend on the fixed profile but not $E$. Hence the trace cost is at most $Ce^{-E}|T|^2/H^2$. Since $H\ge\sigma$, this is a small portion of $(1+\tau)|T|^2/H$ for $E$ sufficiently large in terms of $\sigma$.

All cases retain fixed portions of $A$, $W$, \eqref{radial:h-margin}, and the symmetrized $PG^2$ term, proving \eqref{radial:aligned-margin}. Lemma~\ref{fiber:ranktwo} gives nonpositive real sectional curvature off the fiber axis. The metric is positive there and $h\ge cV$. Smoothness of the original affine fiber potential and its positive limiting metric extend both conclusions across the axis.
\end{proof}
